\scanpage{42}\mspage{4}

\[
u_{(i_2,i_1)(\unsure{i_3},i_4)}
 =u_{(i_1,i_2)(i_3,i_4)}^{-1}\tag*{(2)}
\]

\margin{12 relations reste à fixer 12 quantités}

\[
\begin{aligned}
 &\deleted{\unsure{u}_{\unsure{i_1i_2}}}\\
 &\deleted{u_{(i_2\,i_1)(i_3,\unsure{i})}}
\end{aligned}
\]

\[
\deleted{u_{i_2\,i_1\,i_3\,i_4}}\,
\unsure{u_{(i_1\,i_3)(i_2\,i_4)}}\,
u_{(i_3,i_2)\unsure{(i_1\,i_4)|(i_1\,i_2)}}\,u_{(i_2\,i_1)(i_3\,i_4)}=\mathrm{id}\tag*{(3)}
\]

$V$ est connu en tant que sous-objet de $\prod V_i$, quand on connait le système des $u$, car c'est formé des $(x_i)_{i\in I}$ tels que pour tout $i\in I$ \deleted{\unsure{chaque|donné}} les $x_j$ ($j\in I\setminus\{i\}$) soient ``de somme nulle'' pour le système transitif d'iso entre ces $V_j$ ($j\in J$) défini par l'une ou l'autre (ça revient au $=$) des deux perm. circulaires de $J$.

Mais il faut des conditions pour que le $W$ ainsi défini (qui a priori sera tel que $W\longrightarrow V_i\times V_j=V_J$ \deleted{soit} soit injectif) soit ``assez gros'' pour que ce soit bijectif. On va écrire des conditions pour ceci.

% archive sheet number: 41
